\section{Chapter 6}

\begin{quotex}
\begin{verse}
The energy of middle space the valley is eternal\\
It is the Mysterious Female\\
The gate of the Mysterious Female\\
Is the root of Heaven and Earth\\
Continuous and invariable\\
It acts and does not exhaust itself
\end{verse}
\end{quotex}


There is a possible interaction in this chapter of the doctrine of the Tao with residues of the archaic conception of the primogenial maternal Feminine (\emph{Magna Mater Genitrix,} i.e., Great Mother). ``The Power of middle space'' --- literally in images: ``vital force (or spirit) of the valley''. The valley, understood as the space between two mountain ranges, goes back to the idea both of Emptiness and of the central space where the virtue of the principle is manifested, in the eternal emanation and transformation (contrary to the meaning of the valley --- and of the feminine --- in chapters 61 \& 66).

The gate with two shutters in some commentaries referred to the Principle that acts through the Dyad, the yin and the yang. The whole does not become clear because here the idea of Virtue is superimposed, in a certain way, onto the idea of the Principle. Finally, it has to do with the etheric fluid (chi). This is one of the more convoluted chapters, explained in various ways. It is also about those things that were made objects of esoteric interpretation in operative Taoism, with reference to the human being occultly considered. Characters were used with technical significance whose meaning is difficult to track down. As a hint to interpretations of this type, the Mysterious Female is the principle or energy of the yin soul (\emph{po}) from which those who seek corporeal immortality much extract and coagulate the subtle essence. The gate with two shutters alludes to the nose. The allusion to the technique of breath is learned from the last two lines: ``for long, long stretches one breathes as if one wished to hold the breath, remaining immobile'' (see J J M DeGroot, \emph{Universismus}, Berlin 1918, pp 110-111). The Mysterious Female is the ``interior woman'': ``she is continually in us'', says Kuan Tzu.

\paragraph{Chinese text and literal translation}


Chapter 6 (\begin{CJK*}{UTF8}{bsmi}第六章\end{CJK*})

\columnratio{.4}
\begin{paracol}{2}
\begin{verse}
\begin{CJK*}{UTF8}{bsmi}
谷神不死\\
是謂玄牝\\
玄牝之門\\
是謂天地根\\
綿綿若存\\
用之不勤
\end{CJK*}
\end{verse}
\switchcolumn
\begin{verse}
The valley-spirit will not die,\\
this is the primal mother.\\
The gate of the primal mother\\
is the root of the world.\\
She unceasingly exists,\\
Her influence is without effort.
\end{verse}
\end{paracol}


\flrightit{Posted on 2020-08-29 by Lao Tzu}

